\documentclass[11pt]{article}

\usepackage[a4paper,margin=2cm]{geometry}
\usepackage[T1]{fontenc}
\usepackage{lmodern}
\usepackage{microtype}
\usepackage{amsmath}
\usepackage{booktabs}
\usepackage{enumitem}
\usepackage[hidelinks]{hyperref}
\usepackage{titlesec}

\titlespacing*{\section}{0pt}{8pt}{2pt}

\setlist{nosep,leftmargin=*}
\setlength{\parskip}{4pt}
\setlength{\parindent}{0pt}

\newcommand{\todo}[1]{\textbf{[#1]}}

\begin{document}

\begin{center}
  {\Large\bfseries Project Proposal: Name That LLM}\\[4pt]
  Workshop in Machine Learning and Data Analysis (0368-3548), Fall 2026--27\\[2pt]
  Track: \textbf{competitive} \quad
  Competition: \url{https://www.kaggle.com/competitions/name-that-llm}
\end{center}

\section*{Team}
\begin{tabular}{@{}ll@{}}
  \todo{Full name 1} & \todo{email 1} \\
  \todo{Full name 2} & \todo{email 2} \\
  \todo{Full name 3} & \todo{email 3} \\
\end{tabular}

\section*{The challenge}
Name That LLM is an active Kaggle Featured Code Competition with a \$50{,}000 prize pool
(\$5{,}000 for each of places 1--10). The task is authorship attribution: given a short passage,
predict which of 18 frontier language models wrote it. The 18 models come from six labs:
OpenAI (5 models), Anthropic (5), Google (4), ZAI (2), Qwen (1) and Moonshot AI (1).
Every passage is an authentic, self-contained model response to a prompt on philosophy,
physics, mathematics, biology, computer science or creative prose.

Three properties make the task hard:
\begin{itemize}
  \item \textbf{Unseen prompts.} Test passages answer prompts that do not appear in the
        training set, so the classifier cannot rely on topic.
  \item \textbf{Short, variable length.} Passages are 40--250 words long, and the organisers
        chose lengths so that character counts are not a usable shortcut.
  \item \textbf{Sibling models.} Several models share a lab, training data and post-training
        recipe (e.g.\ five OpenAI models), so the differences between them are subtle.
\end{itemize}

\textbf{Deadlines.} Entry and team merger: 6 January 2027. Final submission: 13 January 2027.
Both fall before 11 April 2027, as the competitive track requires.

\section*{The data}
\begin{itemize}
  \item \texttt{train.csv}: 2{,}699 labelled passages written for 150 prompts, with columns
        \texttt{id}, \texttt{prompt\_id}, \texttt{text} and \texttt{model}. That is roughly
        150 passages per class.
  \item \texttt{test.csv}: about 4{,}000 hidden passages (\texttt{id}, \texttt{text}), visible only
        when the notebook is re-run for scoring. The public leaderboard uses about 40\% of
        them and the private leaderboard the other 60\%; the class mix differs between the two.
  \item Total size of the provided files: 2.14\,MB.
\end{itemize}

\textbf{Metric.} Balanced accuracy, the mean per-class recall over the $K=18$ models:
\[
  \text{BA} = \frac{1}{K}\sum_{k=1}^{K}
  \frac{\left|\{i : y_i = k \wedge \hat{y}_i = k\}\right|}{\left|\{i : y_i = k\}\right|}.
\]
Classes are not equally frequent in the test set, so every model counts the same regardless of
how many passages it has.

\textbf{Submission rules.} Predictions are produced by a Kaggle notebook that runs in at most
9 hours (CPU or GPU) with internet access disabled. Freely and publicly available external
data and pre-trained models are allowed.

\section*{What would count as success}
\begin{enumerate}
  \item A reproducible baseline and a validation scheme we trust: cross-validation grouped by
        \texttt{prompt\_id}, so that validation prompts are unseen, as they are in the test set.
        Success means the cross-validation score tracks the public leaderboard.
  \item A final model that clearly beats our baseline on the private leaderboard. As of
        25 October 2026 the best public score is 0.703 and a score of about 0.64 ranks in the
        top 50 of roughly 3{,}100 teams. We aim for the top 10\% (a bronze medal) and treat a
        higher placement as a stretch goal.
  \item An analysis of what the classifier learns: which models it confuses (we expect
        confusion within a lab), which features separate labs and which separate sibling
        models, and how accuracy changes with passage length.
\end{enumerate}

\section*{Expected compute and storage}
\begin{itemize}
  \item \textbf{Storage.} The data is 2\,MB. The bulk is model weights: a DeBERTa-v3-large
        checkpoint is under 2\,GB, and one open 7--8B LLM for perplexity features is about
        15\,GB in 16-bit precision (about 5\,GB in 4-bit). We expect to stay under 25\,GB, within
        the 30\,GB per-team limit.
  \item \textbf{Compute.} Classical models train in seconds on a CPU. Fine-tuning an encoder on
        2{,}700 short passages takes minutes on a single 16\,GB GPU, and we can run it on
        Kaggle's free GPU quota. The school GPUs are needed mainly for extracting features from
        a 7--8B LLM and for running many cross-validation folds.
  \item \textbf{Inference.} Classifying about 4{,}000 short passages fits easily in the 9-hour
        notebook limit, even with an LLM-based feature extractor.
\end{itemize}

\section*{Initial ideas}
\begin{enumerate}
  \item \textbf{Baseline.} TF-IDF over word and character n-grams with a linear classifier
        (logistic regression or linear SVM), with class weights to match the balanced metric.
  \item \textbf{Stylometric features.} Punctuation habits (em-dashes, semicolons, colons),
        sentence-length statistics, paragraph structure, list and markdown use, and
        characteristic phrases. These are easy to interpret and support the error analysis.
  \item \textbf{Fine-tuned encoders.} DeBERTa-v3 or a similar encoder fine-tuned for 18-way
        classification, with prompt-grouped folds and an ensemble across folds.
  \item \textbf{Hierarchical classification.} Predict the lab first (6 classes, which should be
        easier), then the model within the lab, and compare against flat 18-way classification.
  \item \textbf{Features from open LLMs.} Per-token log-probabilities and perplexity of each
        passage under one or more open LLMs, and hidden-state embeddings as input to a light
        classifier. Models from the same family may score each other's text differently.
  \item \textbf{Ensembling and calibration.} Combine the models above and adjust class
        thresholds for balanced accuracy, since the class mix in the test set differs from
        the training set.
\end{enumerate}

\section*{Timeline and risks}
We plan a working baseline and grouped cross-validation by mid-November, encoder and LLM
feature experiments through December, and the final ensemble before the 13 January deadline.
The main risks:
\begin{itemize}
  \item \textbf{Short schedule.} The final Kaggle deadline is the same day as the midpoint
        review. We will confirm with the staff when the report is due on the competitive track.
        If we miss the Kaggle deadline we will move to the non-competitive track.
  \item \textbf{Small data.} With about 150 examples per class, validation scores are noisy.
        We will rely on grouped cross-validation rather than the public leaderboard.
  \item \textbf{Distribution shift.} The test set has new prompts and a different class mix,
        which can make the public leaderboard misleading.
\end{itemize}

\end{document}
